\section{Listed Futures Daily Settlement}\label{sec:tour-futures}

A listed future behaves unlike the others: it settles every day, it posts
variation margin, and desks run whole engines to manage that margin. This stop shows
that daily settlement needs no engine of its own. It is a scheduled trigger producing
ordinary conserving moves, and the margin that looks like special machinery is just
those moves summing to zero.

The contract is a real one, and it is dollar-denominated, so the tour states it as
such rather than forcing it into euros: ES-FUT, a cash-settled equity-index future,
contract multiplier 50, quoted in USD. Three traders, A, B, and C, and a
clearinghouse hub CH stand behind them. Two rules run the whole life:

\begin{itemize}[leftmargin=1.5em]
\item \textbf{A trade of $\Delta$ contracts at price $p$} adds $-\Delta \cdot p \cdot
50$ to the accumulated cost $ac$ of each side.
\item \textbf{A settle at price $S$} pays each holder variation margin $VM(w) =
\mathit{netq}(w) \cdot S \cdot 50 + ac(w)$, then resets $ac(w) \leftarrow
-\mathit{netq}(w) \cdot S \cdot 50$.
\end{itemize}

The daily settle mark $S$ is not a number the ledger forms. It enters as an externally
sourced observation carrying its basis, and it is consumed under single-basis
consumption: the mark and the positions it is multiplied against must stand in one
frame, exactly as the split stop required.

\subsection{The life of the contract}

Table~\ref{tab:futures-life} is the whole life, in position ($\mathit{netq}$),
accumulated cost ($ac$), and the variation-margin cash each settle pays, for
$(A, B, C)$.

\begin{table}[h]
\centering
\small
\begin{tabular}{@{}llrrr@{}}
\toprule
Step & Event & $\mathit{netq}$ & $ac$ (USD) & VM cash (USD) \\
\midrule
T1        & A buys 10 from B @100 & $(10,-10,0)$ & $(-50{,}000,\ +50{,}000,\ 0)$ & --- \\
Settle d1 & $S_1 = 102$           & $(10,-10,0)$ & $(-51{,}000,\ +51{,}000,\ 0)$ & $(+1{,}000,\ -1{,}000,\ 0)$ \\
T2        & C buys 4 from A @103  & $(6,-10,4)$  & $(-30{,}400,\ +51{,}000,\ -20{,}600)$ & --- \\
Settle d2 & $S_2 = 101$           & $(6,-10,4)$  & $(-30{,}300,\ +50{,}500,\ -20{,}200)$ & $(-100,\ +500,\ -400)$ \\
T3        & B buys 4 from C @101  & $(6,-6,0)$   & $(-30{,}300,\ +30{,}300,\ 0)$ & --- \\
Expiry    & final $S = 105$       & $(6,-6,0)$   & $(-31{,}500,\ +31{,}500,\ 0)$ & $(+1{,}200,\ -1{,}200,\ 0)$ \\
\bottomrule
\end{tabular}
\caption{ES-FUT over its life. Each settle is a scheduled trigger; variation margin is
the cash the settle contract moves.}
\label{tab:futures-life}
\end{table}

Two of the numbers repay a moment. Day 1 is the easy case: A is long 10 into a mark
that rose from 100 to 102, and collects $10 \times 102 \times 50 - 50{,}000 =
\text{USD}~1{,}000$, which B pays. Day 2 looks harder, because A traded intraday. A
sold 4 contracts to C at 103 before the mark fell to 101, and A's day-2 margin is
$6 \times 101 \times 50 - 30{,}400 = -\text{USD}~100$ --- not the naive
$6 \times (101 - 102) \times 50 = -300$ a per-position mark-to-mark would give. The
difference is the \text{USD}~200 A realised selling 4 contracts at 103 into a prior
mark of 102, and the accumulated cost carried it correctly, because $ac$ is a stored,
conserved field and not a re-derivation each day.

\subsection{Margin is just conservation}

The point of the table is the last column. At every settle the variation margin sums
to zero across the holders: $+1{,}000 - 1{,}000 = 0$; $-100 + 500 - 400 = 0$;
$+1{,}200 - 1{,}200 = 0$. This is not a reconciliation the ledger performs after the
fact. It is the same fact as accumulated-cost conservation, $\sum_w \Delta ac = 0$,
which holds because every settle is built from moves and a move is a transfer:
whatever leaves one wallet enters another. The clearinghouse leg is the residual that
balances the holders' legs, and here, with every holder inside the ledger, that
residual is zero each day. Position conservation, $\sum \mathit{netq} = 0$, holds at
every step for the same reason.

Because each settle is keyed by (unit, settlement date), a settle that fires twice ---
a processor retried after a crash, a timer that fired late --- commits once. And the
zero-sum closes over the whole life, not just each day. Cumulative margin is
$A = +2{,}100$, $B = -1{,}700$, $C = -400$, and $2{,}100 - 1{,}700 - 400 = 0$: the
contract's entire economic history is a single conserved stream, checkable end to end.
A's \text{USD}~2{,}100 is exactly $[\,4 \times (103 - 100) + 6 \times (105 - 100)\,]
\times 50$ --- what A made on the 4 contracts sold at 103 and the 6 carried to expiry
at 105 --- so the running margin and the closed-form profit agree to the cent.

One margin lives outside this picture, and the stop should say so plainly. Initial
margin --- the portfolio or SPAN amount a clearinghouse holds against future moves ---
is gross, not zero-sum: every member posts it to the clearinghouse and no member's
posting offsets another's, so it is not a conserved transfer among the holders and does
not belong to the variation-margin stream this table records. It is a collateral
position posted to the clearinghouse, booked like any other collateral through the same
door, and its sizing is a risk model's output the ledger records but never forms. The
zero-sum here is a statement about \emph{variation} margin --- the daily cash the mark
moves between holders --- and about nothing else.

The arithmetic never divides. A settle multiplies $\mathit{netq} \cdot S \cdot m$ and
nothing else --- prices are never added to prices, only multiplied by an integer
position and an integer multiplier --- so there is no place for a rounding error to
enter and accumulate. The ES mark ticks in quarter-points, and a quarter-point on a
multiplier of 50 is an exact $\text{USD}~12.50$, so even the finest price increment
lands on a whole number of cents and the zero-sum holds to the cent at any scale.
\emph{Variation margin is ordinary conserved transactions keyed per settlement day ---
there is no privileged \emph{variation-margin} engine inside the ledger; the same door
and the same log that carry a share purchase carry a cleared future's entire
variation-margin life, while initial margin sits outside as a posted collateral
position.}
